\subsection{Unrestricted smooth critical curves and contracted sections}

\begin{statement}{C.2.UN.9}{Proved}{A smooth full critical curve with generic node is a product node}
Let a reduced hypersurface germ $\{f=0\}$ in a smooth
complex threefold occur in a fibre of a proper torus-surface
degeneration. Suppose that the full reduced critical locus
of $f$ is a smooth curve and its generic transverse
singularity is an ordinary node. There are local
holomorphic coordinates in which
\[
                            f=XY.                  \tag{C.UNNODE.1}
\]
The conclusion about the reduced hypersurface also holds
when the actual map has a holomorphic unit and a common
positive multiplicity.
\end{statement}
\begin{proof}
Choose coordinates $(x,y,z)$ with critical curve
$\Gamma=\{x=y=0\}$. Since $f$ vanishes on $\Gamma$ and
its first derivatives vanish there, its transverse
quadratic term is
\[
                  (x,y)Q(z)(x,y)^{\mathsf t},
\]
where $Q$ is a holomorphic symmetric $2\times2$ matrix.
Its determinant is not identically zero. Shrink the
curve disc so that $\det Q(z)\ne0$ for $z\ne0$.
If $Q(0)$ is invertible, the holomorphic Morse lemma
with parameter $z$ gives $f=X^2+Y^2$, and a linear
change gives (C.UNNODE.1).

Suppose instead that $Q(0)$ is singular. We construct
arbitrarily small constant symmetric matrices $E$
for which $Q(z)+E$ has a rank-one value at a simple
zero of its determinant near the origin. Choose
$z_0\to0$ with $Q'(z_0)\ne0$; such points exist
because the nonconstant determinant makes $Q$
nonconstant. Choose a rank-one symmetric matrix
$R$ close to $Q(0)$ with
\[
                  \operatorname{tr}
                      (\operatorname{adj}(R)Q'(z_0))\ne0.
\]
This choice is possible also when $Q(0)=0$: for a
rank-one symmetric matrix with kernel spanned by
$v$, its adjugate is a nonzero scalar multiple of
$vv^{\mathsf t}$, and the polynomial
$v^{\mathsf t}Q'(z_0)v$ is not identically zero.
Thus the forbidden kernels form a proper finite
set in the projective line, and the rank-one
matrix can be chosen as close to $Q(0)$ as desired.
Set $E=R-Q(z_0)$. Then $E\to0$, and
\[
 \det(Q(z_0)+E)=0,\qquad
 \left.\frac{d}{dz}\det(Q(z)+E)\right|_{z=z_0}\ne0.
 \tag{C.UNNODE.2}
\]

Consider the actual perturbation
$f_E=f+(x,y)E(x,y)^{\mathsf t}$ at $(0,0,z_0)$.
The rank-one Morse lemma, retaining $(y,z)$ as
parameters after a linear normal-coordinate change,
writes
\[
                         f_E=X^2+y^2h(y,z).
\]
The factor $y^2$ follows from $f_E(0,0,z)=0$ and
the vanishing of its normal first derivative along
the same curve. Equation (C.UNNODE.2) says exactly
that $h(0,z_0)=0$ and $\partial h/\partial z(0,z_0)
\ne0$, up to a nowhere-zero coefficient from the
nondegenerate quadratic direction. Hence the actual
coordinate substitution $Z=h(y,z)$ gives
\[
                            f_E=X^2+y^2Z.          \tag{C.UNNODE.3}
\]
The explicit construction in A.5.15 gives a compact
spherical class in a small nonzero level of
(C.UNNODE.3). Its square is
$-2(2/\gcd(2,1))(1/\gcd(2,1))=-4$.
The construction supplies a continuous map from
$S^2$ representing this class; embeddedness is
not required.

Here is the precise transport into an original
fibre. Choose a small closed ball whose boundary
is transverse to $\Gamma$ and to the smooth
stratum of $f=0$. Along the compact boundary link
of $\Gamma$, the matrix $Q$ is invertible.
For all sufficiently small $E$, the parameter
Morse lemma therefore gives product-node charts
there, retaining the curve coordinate. In these
charts $\partial/\partial z$ is tangent to every
level of $f_E$. A real radial curve-direction
vector is transverse to the fixed boundary at
the critical link and remains transverse on its
neighbourhood, uniformly for small $E$. On the
remaining compact boundary, ordinary transversality
gives the same uniform conclusion for small target
values. The full-critical-locus hypothesis ensures
that these charts cover every boundary singular
point requiring this argument.

For a fixed small $E$, transport the spherical
class from (C.UNNODE.3) through regular target values
to a fixed small value $t_1$, avoiding the finitely
many nonzero critical values in the ball. Choose
$E$ small enough that every critical value along
the segment from $E$ to $0$ has modulus less than
$|t_1|/2$. This follows by compactness: a limit of
critical points is a critical point of $f$ with
value zero; the just-proved boundary charts rule
out loss of transversality at the boundary.
The $t_1$-levels along that segment form a proper
submersion with boundary. Ehresmann transport
therefore maps the actual spherical class to a
class of square $-4$ in the original $f=t_1$ level.
Its inclusion into the nearby torus preserves
the square, since both the compact representative
and its intersection push-off lie in that open
level. But this inclusion also maps its representing
$S^2$ to an aspherical torus, where it is
nullhomotopic and has zero homology class. This
contradiction proves $\det Q(0)\ne0$.

If the original map is $u f^m$, replace the reduced
equation by $u^{1/m}f$ using a local holomorphic
root of the unit. Its critical-set germ is still
$\Gamma$: every critical component through the
point has constant value zero, and the singular
locus of the reduced zero hypersurface is unchanged.
The argument then places the
same class in the actual nearby level with base
parameter $t_1^m$. This proves the last assertion.
\end{proof}

\begin{statement}{C.2.UN.10}{Proved}{A contracted irrational horizontal section is elliptic without ambient Betti assumptions}
Let $f:Y\to B$ be a proper holomorphic map from a
smooth complex threefold to a smooth curve, with
connected fibres and smooth general fibre a complex
two-torus. Let $mS$ be an irreducible nonnormal full
fibre. Write $h=\eta\pi:Z\to T\to S$ for the
normalization and its minimal resolution. Suppose
that $Z$ is ruled over a curve of positive genus,
that its effective adjunction divisor is
$D=2C+V$, with $C$ a section and $V$ vertical,
and that $\pi$ contracts $C$. Then the ruling
base has genus one.
\end{statement}
\begin{proof}
We first supply the comparisons that previously
came from the global Betti assumptions. On a
small full fibre tube $U$, the rank-one torsion-free
sheaf $f_*K_U$ is a line bundle. A generator has
no horizontal zeros, since its restriction to a
smooth torus is a nonzero canonical form. Its
zero divisor is consequently $kS$ for some
$k\ge0$, after shrinking away other special
fibres. Thus the evaluation map gives the actual
isomorphism
\[
 K_U\simeq\OO_U(kS),\qquad
 \omega_S\simeq\OO_U((k+1)S)|_S.                \tag{C.UNCR.1}
\]
The divisor $mS$ is principal, so both these line
bundles are torsion. The map to the smooth curve
is flat, and the holomorphic Euler characteristic
of its full fibres is constant. The filtration
of $\OO_{mS}$ has quotients
$\OO_S(-jS)$ for $0\le j<m$; each is a torsion
line-bundle twist of $\OO_S$. Analytic
Riemann--Roch gives the same Euler characteristic
for these twists, since their rational Chern
characters are one. It follows that
\[
            0=\chi(\OO_{mS})=m\,\chi(\OO_S).
 \tag{C.UNCR.2}
\]
The exact normalization and resolution duality
maps therefore give $K_Z+D=h^*\omega_S$, with
the same effective divisor and torsion right-hand
side as in C.2.CR.1--C.2.CR.2.

Here are the parts of that argument needed here,
with no transgression-class assumption. Finite
duality embeds $\omega_T$ into the torsion line
$\eta^*\omega_S$, with its sections vanishing on
the nonempty conductor. A nonzero section of a
torsion line on a connected compact normal space
is nowhere zero: taking a trivializing power
makes it a nonzero constant. Consequently
$H^2(T,\OO_T)=0$. The coherent Leray sequence
then bounds $\operatorname{length}R^1\pi_*\OO_Z$
by $h^1(Z,\OO_Z)=g$, whereas formal functions
and restriction to the contracted section $C$
give length at least $h^1(C,\OO_C)=g$ at its
image point $p$. Thus equality holds at $p$ and
there are no further contributions.

At $q=\eta(p)$, fix a local generator of $\omega_S$.
Finite duality identifies the branch conductor
ideal $I_p$ with $\omega_{T,p}$, embedded as an
actual ideal summand in $\OO_{S,q}$. Proper
duality and the residue sequence of $C$ give
exactly
\[
 I_p=(\pi_*\omega_Z(C))_p,\qquad
 I_p/(\pi_*\omega_Z)_p\simeq H^0(C,\omega_C).
 \tag{C.UNCR.3}
\]
The equality of lengths follows from the preceding
formal-functions computation. In the fixed local
dualizing trivialization these sheaves are
$\OO_Z(-C-V)$ and $\OO_Z(-2C-V)$.
The maximal ideal of $\OO_{S,q}$ vanishes on $C$;
its square supplies an additional order along
$C$ without losing any prescribed order along
$V$. Hence
$I_p\cap\mathfrak m_q^2\subset(\pi_*\omega_Z)_p$.
The embedding dimension three and Nakayama
argument of C.2.CR.2 now give $g\le2$. When
$g=2$ they give the actual conductor ideal
$(x,y)$, a unique normalization point $p$ above
$q$, and generic transverse delta invariant
one. Thus the full conductor-image curve is
smooth and reduced at $q$, and the generic
transversal is a node or a cusp.

Every positive-dimensional critical component
of the reduced surface is part of that conductor
curve; after shrinking, its full reduced critical
locus is exactly the smooth curve just obtained.
The node case is a product node by C.2.UN.9.
Its normalization has two smooth points above
$q$, contradicting the unique point $p$.
The cusp case is a product cusp by C.2.CR.4.
Its explicit normalization
$(u,z)\mapsto(X,Y,z)=(u^3,-u^2,z)$, after
adjusting the harmless sign in $X^2+Y^3$,
is smooth. This contradicts a minimal
normalization resolution contracting a
positive-genus curve to $p$. Both genus-two
possibilities are excluded, and $g=1$ follows.
\end{proof}
